%% ============================================================
%% thesis_outline.tex
%% Bachelor Thesis — Vrije Universiteit Amsterdam
%% LLM Agents for Interlinked and Personalised Knowledge Bases
%% Student: Muhammed Furkan Kaya
%% Supervisor: Emile van Krieken
%% ============================================================

\documentclass[12pt, a4paper]{article}

%% ── Packages ─────────────────────────────────────────────────
\usepackage[utf8]{inputenc}
\usepackage[T1]{fontenc}
\usepackage[english]{babel}
\usepackage{geometry}
\geometry{top=2.5cm, bottom=2.5cm, left=3cm, right=3cm}

\usepackage{setspace}
\onehalfspacing

\usepackage{hyperref}
\hypersetup{colorlinks=true, linkcolor=blue, citecolor=blue, urlcolor=blue}

\usepackage{amsmath, amssymb}
\usepackage{graphicx}
\usepackage{booktabs}
\usepackage{multirow}
\usepackage{xcolor}
\usepackage{listings}
\usepackage{caption}
\usepackage{subcaption}
\usepackage{enumitem}
\usepackage{cleveref}
\usepackage{natbib}  % or biblatex — switch to what VU requires

%% ── Title ────────────────────────────────────────────────────
\title{
    \textbf{LLM Agents for Interlinked and Personalised Knowledge Bases} \\
    \large Bachelor Thesis \\
    \normalsize Vrije Universiteit Amsterdam \\
    Programme: B Artificial Intelligence
}
\author{
    Muhammed Furkan Kaya \\
    \small Student ID: 2780604 \\
    \small \texttt{m.f.kaya@student.vu.nl} \\[1em]
    \small Supervisor: Emile van Krieken \\
    \small Second assessor: Frank van Harmelen \\
    \small Department of Computer Science, AI: Learning and Reasoning
}
\date{Academic Year 2025--2026}

%% ============================================================
\begin{document}
\maketitle
\thispagestyle{empty}

%% ── Abstract ─────────────────────────────────────────────────
\begin{abstract}
% TBD after experiments. Draft placeholder:
Personal knowledge management tools such as Obsidian have become central to academic
knowledge work. However, maintaining a structured, interlinked knowledge base over a
growing corpus of academic papers requires significant manual effort. This thesis
investigates whether a large language model (LLM) agent pipeline can automate the
creation of structured, interlinked Obsidian-style notes from academic PDFs and,
crucially, whether this structured representation improves retrieval quality compared
to a flat chunk-based retrieval-augmented generation (RAG) baseline. We construct a
curated corpus of 15--25 academic papers and evaluate retrieval quality using
Recall@k, Precision@k, F1@k, and Mean Reciprocal Rank (MRR). We further evaluate
the quality of agent-generated metadata and inter-note links using field-level
accuracy and link Precision/Recall/F1. A supplementary expert evaluation assesses
summary quality on a sample of 20 notes. Our results show [TBD]. We discuss
implications for LLM-based personal knowledge management and the limitations of
the proposed approach.

\noindent \textbf{Keywords:} retrieval-augmented generation, knowledge graph,
personal knowledge management, Obsidian, LLM agents, structured notes, metadata
extraction, inter-note linking, information retrieval.
\end{abstract}

\tableofcontents
\newpage

%% ============================================================
\section{Introduction}
\label{sec:intro}

%% Motivation: PKM tools + growing need for structure
%% Obsidian as a knowledge graph (cite Emile's blog / template)
%% The manual overhead problem
%% LLM agents as a potential solution

\subsection{Context and Motivation}
% - Academics face information overload.
% - PKM tools (Obsidian) treat vaults as knowledge graphs.
% - Manual structuring (metadata, links) is time-consuming.
% - LLM agents can automate text understanding and structured output.

\subsection{Problem Statement}
% A controlled computational study comparing two systems:
% (1) Baseline RAG: flat chunk-based retrieval
% (2) Proposed: LLM-agent-generated structured Obsidian notes with metadata + links
% Core question: does the structure help retrieval?

\subsection{Research Questions}
\label{sec:rqs}

\noindent \textbf{Main Research Question:}

\begin{quote}
To what extent can an LLM agent improve retrieval in an Obsidian-style personal
knowledge base by generating structured metadata and inter-note links from academic
papers, compared to a simple RAG baseline?
\end{quote}

\noindent \textbf{Sub-questions:}

\begin{itemize}
    \item[\textbf{RQ1}] \textit{Retrieval performance}: Does the structured,
        interlinked note system improve retrieval quality (Recall@k, Precision@k,
        F1@k, MRR) compared to a flat RAG baseline?

    \item[\textbf{RQ2}] \textit{Metadata quality}: How accurately does the LLM
        agent generate metadata fields (title, authors, year, venue, topics,
        methods, datasets, evaluation metrics, contributions)?

    \item[\textbf{RQ3}] \textit{Link quality}: How accurately and usefully does
        the LLM agent generate inter-note links, and how correct are the assigned
        relation types?

    \item[\textbf{RQ4}] \textit{Summary quality (supplementary)}: How does the
        supervisor rate the factual correctness and usefulness of generated note
        summaries on a sample of 20 notes?
\end{itemize}

\subsection{Contributions}
% 1. A pipeline for LLM-agent-based structured Obsidian note generation from PDFs
% 2. A controlled comparison of flat RAG vs. structured notes for retrieval
% 3. Evaluation metrics and annotation guidelines for metadata/link quality
% 4. Reusable query set, gold labels, code, and prompts (open science)

\subsection{Thesis Outline}
% Point to each section.

%% ============================================================
\section{Related Work}
\label{sec:related}

\subsection{Retrieval-Augmented Generation}
% Lewis et al. 2020 — foundational RAG
% Fusion-in-Decoder (Izacard & Grave 2020)
% Focus: flat chunking, dense retrieval, metrics used

\subsection{Graph-Structured and Knowledge-Graph-Enhanced RAG}
% GraphRAG (Edge et al. 2024) — community summaries, global queries
% LightRAG (Guo et al. 2024) — dual-level retrieval
% HippoRAG (Gutierrez et al. 2024) — hippocampal KG, Recall@k evaluation
% RAPTOR (Sarthi et al. 2024) — hierarchical tree retrieval

\subsection{Agentic Memory and Structured Memory Systems}
% A-MEM (Weng et al. 2025) — Zettelkasten-inspired memory notes + links
% MemGPT (Packer et al. 2023) — OS-inspired agent memory
% MemoryBank (Zhong et al. 2023) — long-term LLM memory
% MemoRAG (Qian et al. 2024) — memory clues for fuzzy queries

\subsection{Personal Knowledge Graphs}
% Balog & Kenter 2019 — PKG research agenda, link quality challenges
% Connection to Obsidian as a personal knowledge graph

\subsection{Evaluation Frameworks}
% RAGAS (Es et al. 2023) — automated RAG evaluation, why I avoid LLM-as-judge
% LoCoMo (Maharana et al. 2024) — gold label annotation design, stratified evaluation

\subsection{LLM-Agent Pipelines for Structured Content}
% STORM (Shao et al. 2024) — multi-agent pipeline for structured writing
% My Brain Is Full Crew — practitioner inspiration (design reference, not formal baseline)
% Emile's academic-obsidian template — target note schema design

\subsection{Gap Analysis}
% Table: prior work vs. this thesis
% Key gap: no prior work evaluates typed link quality in structured personal academic notes
% with standard IR metrics against a flat RAG baseline.

%% ============================================================
\section{Methodology}
\label{sec:methodology}

\subsection{Research Design}
% Controlled computational experiment.
% No user study. No human participants. No consent forms.
% Independent variable: retrieval system (Baseline RAG vs. Structured Notes)
% Dependent variable: retrieval quality metrics
% Held constant: corpus, query set, embedding model

\subsection{Dataset Construction}
\label{sec:dataset}

\subsubsection{Paper Selection}
% Purposive selection from arXiv and public academic sources.
% Selection criteria:
%   - Publicly available (arXiv or open access)
%   - Relevant to thesis topic categories (RAG, GraphRAG, PKG, agentic memory, etc.)
%   - Non-trivial semantic diversity (similar but not identical topics)
% Target: 15--25 papers
% See Table~\ref{tab:corpus} for selected papers.

\subsubsection{Text Extraction}
% PyMuPDF (fitz) text extraction.
% Cleaning: dehyphenation, whitespace normalisation.
% Exclusion if extraction quality is poor.

\subsubsection{Query Set Design}
\label{sec:querydesign}
% 50--100 queries, manually constructed and reviewed.
% Stratified by difficulty type:

\begin{table}[h]
\centering
\caption{Query difficulty types and example queries.}
\label{tab:querytypes}
\begin{tabular}{@{}lll@{}}
\toprule
\textbf{Type} & \textbf{Description} & \textbf{Example} \\
\midrule
Simple factual & Single-paper lookup & ``Which paper introduced RAG?'' \\
Method query   & Single-paper method & ``What retrieval method does GraphRAG use?'' \\
Comparison     & Cross-paper relationship & ``How does GraphRAG differ from standard RAG?'' \\
Cross-note     & Multi-paper synthesis & ``Which papers argue structure improves retrieval?'' \\
Metadata       & Field-specific lookup & ``Which papers evaluate using Recall@k?'' \\
Link/relation  & Relationship query & ``Which papers are related to agentic memory?'' \\
Failure-prone  & Semantic ambiguity & ``Which systems use graph structure for reasoning?'' \\
\bottomrule
\end{tabular}
\end{table}

% ~50\% simple/single-note, ~50\% complex/cross-note queries.
% Gold relevant documents: manually annotated binary relevance per (query, note) pair.
% Annotation process: author judges relevance; spot-checked by supervisor.

\subsection{Systems Compared}
\label{sec:systems}

\subsubsection{Baseline RAG System}
% Input: same academic PDFs.
% Processing:
%   (1) Extract text (PyMuPDF)
%   (2) Fixed-size chunking: 800 tokens, 120-token overlap
%   (3) Embed chunks: all-MiniLM-L6-v2 (sentence-transformers)
%   (4) Store in NumPy cosine similarity index
%   (5) Retrieve top-k chunks for each query
% Output: ranked list of chunks mapped to source paper IDs.
% No generated metadata. No inter-document links.

\subsubsection{Proposed: LLM-Agent Structured Note System}
% Input: same academic PDFs.
% Agent pipeline:
%   (1) PDF Ingestion Agent: extract text
%   (2) Metadata Agent: LLM extracts structured metadata (see note schema)
%   (3) Paper Note Agent: LLM generates Obsidian Markdown note
%   (4) Link Agent: for each note, retrieve candidate notes via embedding similarity,
%       LLM classifies link type and direction
% Retrieval:
%   (5) Embed full notes (frontmatter + content concatenated)
%   (6) Retrieve top-k notes by cosine similarity
% Output: ranked list of notes.

% Note schema:
\begin{verbatim}
---
type: source/paper
title: "..."
aliases: [...]
authors: [...]
year: 2024
venue: "..."
doi_or_arxiv: "..."
hasTopic: [...]
hasConcept: [...]
usesMethod: [...]
usesDataset: [...]
evaluatesWith: [...]
relatedPapers: []
links_to: []
created_by: "LLM-agent"
---
\end{verbatim}

% Link types: hasTopic, usesMethod, evaluatesOn, extends, comparesTo, contradicts,
%             hasDataset, publishedIn, isA, subset.

\subsubsection{Ablation: Link-Expanded Retrieval}
% Step 1: retrieve top-k notes (same as proposed system).
% Step 2: add all notes directly linked to top-k results.
% Re-rank combined set by embedding similarity score.
% Evaluates whether link expansion improves Recall@k at the cost of Precision@k.

\subsection{LLM Configuration}
% Provider: [TBD — Gemini 1.5 Flash / GPT-4o-mini / Claude 3 Haiku]
% Temperature: 0.0 for metadata and link generation; 0.3 for note/summary generation
% Max output tokens: 2048 for metadata; 4096 for notes
% Prompts: stored verbatim in prompts/ directory (see Appendix~\ref{app:prompts})
% Model version pinned and documented in configs/structured.yaml

\subsection{Evaluation Data}
\label{sec:evaldata}

\subsubsection{Gold Relevance Labels for Retrieval}
% Binary: doc\_id is relevant to query (1) or not (0).
% Annotated by author for all (query, paper) pairs.
% Spot-checked by supervisor for a random sample of 20 pairs.

\subsubsection{Gold Metadata Labels}
% Manually annotated for each paper:
%   - Exact-match fields: title, year, venue, doi\_or\_arxiv
%   - Multi-label fields: topics, methods, datasets, metrics, concepts
% Source: paper itself (first page, abstract, experiment section)
% Annotation time: ~15--20 minutes per paper

\subsubsection{Gold Link Labels}
% For each pair of papers in the corpus:
%   - Is there a meaningful link? (binary)
%   - If yes, what link type? (hasTopic, usesMethod, extends, comparesTo, etc.)
%   - Annotation: author constructs gold link set based on reading both papers.
%   - Spot-checked by supervisor.

\subsubsection{Summary Evaluation Sample}
% 20 generated summaries randomly sampled from the corpus.
% Supervisor rates each on a 1--5 scale on five dimensions:
%   (1) Factual correctness
%   (2) Coverage of main contribution
%   (3) Absence of hallucination
%   (4) Usefulness for rediscovery
%   (5) Clarity
% Not used as a primary thesis result; reported as supplementary.

%% ============================================================
\section{Evaluation Metrics}
\label{sec:metrics}

\subsection{Retrieval Quality Metrics}

Let $Q$ be the set of queries, $R_q$ the ranked list of retrieved documents for query $q$,
and $G_q$ the set of gold relevant documents.

\textbf{Recall@k:}
\begin{equation}
\text{Recall@}k = \frac{|R_q^{(k)} \cap G_q|}{|G_q|}
\end{equation}
where $R_q^{(k)}$ denotes the top-$k$ retrieved documents. Measures coverage of relevant documents.

\textbf{Precision@k:}
\begin{equation}
\text{Precision@}k = \frac{|R_q^{(k)} \cap G_q|}{k}
\end{equation}
Measures the fraction of retrieved documents that are relevant.

\textbf{F1@k:}
\begin{equation}
\text{F1@}k = \frac{2 \cdot \text{Precision@}k \cdot \text{Recall@}k}{\text{Precision@}k + \text{Recall@}k}
\end{equation}

\textbf{Mean Reciprocal Rank (MRR):}
\begin{equation}
\text{MRR} = \frac{1}{|Q|} \sum_{q \in Q} \frac{1}{\text{rank}(q)}
\end{equation}
where $\text{rank}(q)$ is the rank of the first relevant document. Emphasises top-ranking quality.

All metrics computed for $k \in \{1, 3, 5, 10\}$ and macro-averaged over all queries.
Results also reported stratified by query type (simple/single-note vs complex/cross-note).

\subsection{Metadata Quality Metrics}

\textbf{Exact match} for scalar fields:
\begin{equation}
\text{ExactMatch}(\hat{f}, f^*) = \mathbf{1}[\hat{f} = f^*]
\end{equation}
Applied to: \texttt{year}, \texttt{venue}, \texttt{doi\_or\_arxiv} (after normalisation).

\textbf{Precision, Recall, F1} for multi-label fields:
\begin{align}
\text{Precision}_f &= \frac{|\hat{F} \cap F^*|}{|\hat{F}|} \\
\text{Recall}_f &= \frac{|\hat{F} \cap F^*|}{|F^*|} \\
\text{F1}_f &= \frac{2 \cdot \text{Precision}_f \cdot \text{Recall}_f}{\text{Precision}_f + \text{Recall}_f}
\end{align}
Applied to: \texttt{topics}, \texttt{concepts}, \texttt{keywords}, \texttt{evaluation\_metrics},
\texttt{datasets}, \texttt{method}, \texttt{main\_contributions}.

All metrics macro-averaged across the corpus.

\subsection{Link Quality Metrics}

Let $\hat{L}$ be the set of generated links and $L^*$ the gold link set.

\begin{align}
\text{Link Precision} &= \frac{|\hat{L} \cap L^*|}{|\hat{L}|} \\
\text{Link Recall} &= \frac{|\hat{L} \cap L^*|}{|L^*|} \\
\text{Link F1} &= \frac{2 \cdot \text{LP} \cdot \text{LR}}{\text{LP} + \text{LR}}
\end{align}

A generated link $\hat{l} = (\text{source}, \text{target}, \text{type})$ is counted as
correct if both the pair and the type match a gold link. If only the pair matches (type
incorrect), it counts as a partial match and is reported separately.

Additional reported counts:
\begin{itemize}
    \item Hallucinated links: $|\hat{L} \setminus L^*|$ (generated but not in gold)
    \item Missing links: $|L^* \setminus \hat{L}|$ (in gold but not generated)
    \item Relation-type accuracy: among correctly identified pairs, fraction with correct type
\end{itemize}

\subsection{Summary Quality (Supplementary)}
% 1--5 Likert scale on 5 dimensions, rated by supervisor.
% Mean score per dimension reported.
% No statistical significance testing (n=20, supplementary only).

%% ============================================================
\section{Data Management and Open Science}
\label{sec:datamanagement}

\subsection{Data Sources}
% - Public arXiv PDFs (not redistributed; arXiv IDs shared)
% - LLM-generated notes, metadata, links (fully shareable)
% - Query set and gold labels (fully shareable)
% - Evaluation outputs and metric results (fully shareable)

\subsection{No Personal Data}
% No user study. No participant data. No consent forms required.
% No private Obsidian vaults, emails, calendars, or personal notes.
% All data is derived from public academic papers.

\subsection{LLM Usage Transparency}
% Provider, model name, and version pinned and reported.
% Temperature and generation parameters documented in configs/.
% Prompts published verbatim in GitHub repository.

\subsection{Reproducibility}
% GitHub repository: [URL TBD]
% Contents: code, prompts, configs, query set, gold labels, evaluation scripts.
% PDFs not included (copyright). arXiv IDs provided in data/raw\_pdfs/manifest.json.
% Random seeds documented where applicable.
% Download script provided for corpus reconstruction.

\subsection{Licensing}
% Code: MIT license.
% Emile's academic-obsidian template used as schema reference only (GPL-3.0).
% Note schema and prompt design inspired by the template; no files copied.

%% ============================================================
\section{Planned Experiments}
\label{sec:experiments}

\subsection{Main Comparison}
% Baseline RAG vs Structured Notes on 50--100 queries.
% Primary metric: Recall@5, F1@5, MRR.
% Expected finding: structured notes improve cross-note queries more than simple queries.

\subsection{Ablation Studies}
\begin{enumerate}
    \item Flat RAG vs Structured Notes (no links) vs Structured Notes + Link Expansion
    \item Metadata included in indexed text vs excluded
    \item Simple queries vs complex/cross-note queries (stratified analysis)
\end{enumerate}

\subsection{Metadata Quality Evaluation}
% Run Metadata Agent on all papers.
% Compare against gold labels.
% Report per-field accuracy.

\subsection{Link Quality Evaluation}
% Run Link Agent on all notes.
% Compare against gold link set.
% Report Link P/R/F1 and relation-type accuracy.

\subsection{Summary Quality (Supplementary)}
% Sample 20 notes. Share summaries with supervisor.
% Supervisor completes rubric. Report mean ratings.

%% ============================================================
\section{Results}
\label{sec:results}
% TBD after experiments.

%% ============================================================
\section{Discussion}
\label{sec:discussion}

\subsection{Main Findings}
% Interpret results w.r.t. RQ1--RQ4.
% When does structure help? When does it not?

\subsection{Analysis by Query Type}
% Simple vs complex query stratification.
% Hypothesis: structure helps more for cross-note queries.

\subsection{Failure Analysis}
% Cases where structured system is worse.
% Hallucinated links. Missing links. Wrong metadata.

\subsection{Limitations}
\begin{itemize}
    \item Small corpus (15--25 papers). Results may not generalise.
    \item Purposive sampling may bias results.
    \item Query set designed by same author who built the system.
    \item Gold labels involve human judgment.
    \item LLM outputs vary across versions; results tied to specific model + prompts.
    \item Results depend on embedding model, chunking parameters.
    \item Controlled academic vault ≠ messy real-world personal vault.
    \item Summary quality evaluation is supplementary and based on 20 samples.
\end{itemize}

\subsection{Future Work}
% User study with real Obsidian users.
% Larger corpora.
% Graph-based retrieval (not just vector).
% Multi-vault or cross-user personalisation.

%% ============================================================
\section{Conclusion}
\label{sec:conclusion}
% Summarise thesis in 3--4 paragraphs.
% Answer main RQ concisely.
% Emphasise what is new (typed link quality evaluation, structured note retrieval, etc.)

%% ============================================================
%% References
%% ============================================================
\bibliographystyle{plainnat}
\bibliography{refs}

%% ============================================================
%% Appendices
%% ============================================================
\appendix

\section{Paper Corpus}
\label{app:corpus}
% Table of all 15--25 papers with arXiv ID, title, year, category.

\section{Prompts}
\label{app:prompts}
% Full text of metadata_agent.md, paper_note_agent.md, link_agent.md, summary_agent.md.

\section{Query Set}
\label{app:queries}
% Full list of 50--100 queries with difficulty type and gold labels.

\section{Annotation Guidelines}
\label{app:annotation}
% Instructions for gold label construction (metadata, links, relevance).
% Based on docs/annotation_guidelines.md.

\section{Summary Evaluation Rubric}
\label{app:rubric}
% The 5-dimension rubric used for supervisor summary review.

\end{document}
